% Part: normal-modal-logic
% Chapter: axioms-systems
% Section: normal-logics

\documentclass[../../../include/open-logic-section]{subfiles}

\begin{document}

\olfileid{nml}{prf}{nor}

\olsection{منطق‌های موجهات عادی}

هر مجموعه‌ای از !!{formula}های موجهاتی را نمی‌توان به‌آسانی به‌صورت
!!{formula}های اشتقاق‌پذیر از مجموعه‌ای از اصل‌ها مشخص کرد. می‌خواهیم
منطق‌های موجهات خوش‌رفتار باشند. نخست، هرچه را در منطق گزاره‌ای کلاسیک
می‌توانیم !!{derive} کنیم باید همچنان !!{derivable} باشد؛ البته با توجه
به اینکه !!{formula}ها اکنون ممکن است
\iftag{prvBox}{$\Box$\iftag{prvDiamond}{
    و~}{}}{}\iftag{prvDiamond}{$\Diamond$}{} را نیز در بر داشته باشند.
برای این منظور لازم می‌دانیم یک منطق موجهات شامل همهٔ نمونه‌های
همان‌گویی باشد و تحت وضع مقدم بسته باشد.

\begin{defn}
  یک \emph{منطق موجهات} مجموعهٔ~$\Sigma$ از !!{formula}های موجهاتی است
  که
  \begin{enumerate}
  \item شامل همهٔ همان‌گویی‌هاست؛ و
  \item تحت جانشینی بسته است؛ یعنی اگر~$!A \in \Sigma$ و~$!D_1$،
    \dots، $!D_n$ !!{formula} باشند، آنگاه
    \[
    \SSubst{!A}{\subst{!D_1}{p_1}, \dots, \subst{!D_n}{p_n}} \in \Sigma
    \]؛
  \item تحت \emph{وضع مقدم} بسته است؛ یعنی اگر~$!A$ و
    $!A \lif !B \in \Sigma$، آنگاه~$!B \in \Sigma$.
  \end{enumerate}
\end{defn}

برای استفاده از معناشناسی رابطه‌ای برای منطق‌های موجهات، باید همچنین
همهٔ !!{formula}های معتبر در همهٔ مدل‌های موجهاتی را در آن‌ها بگنجانیم.
روشن می‌شود این شرط به‌محض آن برآورده است که همهٔ نمونه‌های
\Ax{K}\iftag{prvDiamond}{ و~\Dual{}}{} !!{derivable} باشند و هرگاه
!!a{formula}~$!A$ !!{derivable} است، $\Box !A$ نیز چنین باشد. منطق
موجهاتی که این شرایط را برآورده کند \emph{عادی} نامیده می‌شود. (البته
منطق‌های موجهات غیرعادی نیز وجود دارند، اما مدل‌های رابطه‌ای معمول برای
آن‌ها مناسب نیستند.)

\begin{defn}
  منطق موجهات~$\Sigma$ \emph{عادی} است اگر شامل
  \begin{align*}
    \tag{\Ax{K}} & \Box(p \lif q) \lif (\Box p \lif \Box q),
    \iftag{prvDiamond}{\\
      \tag{\Dual} & \Diamond p \liff \lnot\Box\lnot p}{}
  \end{align*}
  باشد و تحت \emph{ضرورت‌بخشی} بسته باشد؛ یعنی اگر~$!A \in \Sigma$،
  آنگاه~$\Box !A \in \Sigma$.
\end{defn}

توجه کنید که درحالی‌که استلزام همان‌گویانه به‌اندازهٔ کافی «ریزساخت»
است تا \emph{صدق در یک جهان} را حفظ کند، قاعدهٔ~\Nec{} تنها \emph{صدق
در یک مدل} را حفظ می‌کند (و ازاین‌رو اعتبار در یک قاب یا رده‌ای از
قاب‌ها را نیز حفظ می‌کند).

\begin{prop}\ollabel{prop:rk}
  هر منطق موجهات عادی تحت قاعدهٔ~\RK بسته است:
  \begin{prooftree}
    \AxiomC{$!A_1 \lif (!A_2 \lif \cdots (!A_{n-1} \lif !A_n)\cdots)$}
    \RightLabel{\RK}
    \UnaryInfC{$\Box!A_1 \lif (\Box!A_2 \lif \cdots (\Box!A_{n-1}
      \lif \Box!A_n)\cdots).$}
  \end{prooftree}
\end{prop}

\begin{proof}
  با استقرا بر~$n$: اگر~$n=1$، قاعده همان~\Nec{} است و هر منطق
  موجهات عادی تحت~\Nec{} بسته است.

  اکنون فرض کنید نتیجه برای~$n-1$ برقرار باشد؛ نشان می‌دهیم برای~$n$
  نیز برقرار است.

  فرض کنید
  \begin{align*}
  & !A_1 \lif (!A_2 \lif \cdots (!A_{n-1} \lif !A_n)\cdots) \in \Sigma
  \intertext{بنا بر فرض استقرا داریم}
  & \Box!A_1 \lif (\Box!A_2 \lif \cdots \Box(!A_{n-1} \lif !A_n)\cdots)
  \in \Sigma
  \intertext{چون~$\Sigma$ منطق موجهات عادی است، همهٔ نمونه‌های~$\Ax{K}$،
    و به‌ویژه عبارت زیر، را شامل می‌شود:}
  & \Box(!A_{n-1} \lif !A_n) \lif (\Box!A_{n-1} \lif \Box!A_n) \in \Sigma
  \intertext{با استفاده از وضع مقدم و نمونه‌های همان‌گویانهٔ مناسب به دست می‌آوریم}
  & \Box!A_1 \lif (\Box!A_2 \lif \cdots (\Box!A_{n-1}
  \lif \Box!A_n)\cdots) \in \Sigma. 
  \end{align*}
\end{proof}

\begin{prop}\ollabel{prop:notDiamondBot}
  هر منطق موجهات عادی~$\Sigma$ شامل~$\lnot\Diamond\lfalse$ است.
\end{prop}

\begin{prob}
  \olref[nml][prf][nor]{prop:notDiamondBot} را اثبات کنید.
\end{prob}

\begin{prop}
  فرض کنید~$!A_1$، \dots، $!A_n$ !!{formula} باشند. آنگاه کوچک‌ترین
  منطق موجهات عادی~$\Sigma$ که همهٔ نمونه‌های~$!A_1$، \dots، $!A_n$
  را در بر دارد وجود دارد.
\end{prop}

\begin{proof}
  با داشتن~$!A_1$، \dots، $!A_n$، $\Sigma$ را اشتراک همهٔ منطق‌های
  موجهات عادی‌ای تعریف کنید که همهٔ نمونه‌های~$!A_1$، \dots، $!A_n$
  را در بر دارند. این خانواده ناتهی است، زیرا~$\Frm[L]$، یعنی مجموعهٔ
  همهٔ !!{formula}ها، یک چنین منطق موجهاتی است.
\end{proof}

\begin{defn}
کوچک‌ترین منطق موجهات عادی شامل~$!A_1$، \dots، $!A_n$ یک
\emph{دستگاه موجهاتی} نامیده می‌شود و با~$\Log{K} !A_1 \dots !A_n$
نشان داده می‌شود. کوچک‌ترین منطق موجهات عادی با~\Log{K} نشان داده می‌شود.
\end{defn}

\end{document}
